% -----------------------------------------------------------------------------------------------------
\section{多腿振幅的红外极点}\label{sec:sudakov-multileg}

形状因子只有两条有色外线，其色结构是平凡的。$n$ 条有色外线的在壳振幅在维数正规化中同样有至多 $1/\varepsilon^{2L}$ 的红外极点，而且由于软胶子可以连接任意两条外线，
极点带有非平凡的色关联。描述它的自然语言是\emph{色空间形式}：振幅是色空间中的矢量 $|\mathcal M(\{p\},\varepsilon)\rangle$，$\mathbf T_i$ 是作用于第 $i$ 个部分子色指标的生成元
（出射夸克为 $t^a$，出射反夸克为 $-t^{a\,\tp}$，胶子为 $\ii f^{\cdot a\cdot}$ 型，入射粒子取共轭），$\mathbf T_i\cdot\mathbf T_j\equiv\mathbf T_i^a\mathbf T_j^a$，$\mathbf T_i^2=C_i$（$C_q=C_F$，$C_g=C_A$），
色守恒为 $\sum_i\mathbf T_i|\mathcal M\rangle=0$。本节说明，形状因子的全部结构——因子化、cusp 反常维度、指数化——推广为一个色空间中的乘性重整化因子 $\mathbf Z$，
而两圈以内它的反常维度只含两两的“偶极子”色关联。

\begin{theorem}[无质量多腿振幅红外极点的因子化 \tagA]\label{thm:sudakov-ir-poles}
设 $|\mathcal M(\{p\},\varepsilon)\rangle$ 是 $n$ 个无质量部分子的紫外已重整化的在壳振幅（红外由维数正规化正规化）。则存在色空间中的矩阵 $\mathbf Z(\{p\},\mu,\varepsilon)$（只含 $\varepsilon$ 的极点），
使 $|\mathcal M_{\rm fin}(\{p\},\mu)\rangle\equiv\mathbf Z^{-1}|\mathcal M(\{p\},\varepsilon)\rangle$ 在 $\varepsilon\to0$ 时有限：
\begin{equation}\label{eq:sudakov-Zfact}
|\mathcal M(\{p\},\varepsilon)\rangle=\mathbf Z(\{p\},\mu,\varepsilon)\,|\mathcal M_{\rm fin}(\{p\},\mu)\rangle,\qquad
\frac{\dd\mathbf Z}{\dd\ln\mu}=-\boldsymbol\Gamma\,\mathbf Z,\qquad \frac{\dd}{\dd\ln\mu}|\mathcal M_{\rm fin}\rangle=\boldsymbol\Gamma\,|\mathcal M_{\rm fin}\rangle .
\end{equation}
两圈以内，反常维度取偶极子形式
\begin{equation}\label{eq:sudakov-dipole}
\boldsymbol\Gamma=\sum_{(i,j)}\frac{\mathbf T_i\cdot\mathbf T_j}{2}\,\gamma_{\rm cusp}(\alpha_s)\ln\frac{\mu^2}{-s_{ij}}+\sum_i\gamma^i(\alpha_s),\qquad s_{ij}=2\sigma_{ij}\,p_i\cdot p_j+\ii0,
% VF-SUD-5
\end{equation}
其中 $(i,j)$ 遍历 $i\ne j$ 的全部\emph{有序}对（每个无序对计两次），$\sigma_{ij}=+1$（$i,j$ 同为入射或同为出射）或 $-1$（其余），
$\gamma_{\rm cusp}=4\frac{\alpha_s}{4\pi}+O(\alpha_s^2)$ 是去掉 Casimir 的 cusp 反常维度（$\Gamma^F_{\rm cusp}=C_F\gamma_{\rm cusp}$，\cref{def:sudakov-cusp}），$\gamma^q=-3C_F\frac{\alpha_s}{4\pi}+O(\alpha_s^2)$。
在单圈阶，$\ln\mathbf Z=\frac{\alpha_s}{4\pi}\big[\frac{\boldsymbol\Gamma'_0}{4\varepsilon^2}+\frac{\boldsymbol\Gamma_0}{2\varepsilon}\big]+O(\alpha_s^2)$，其中 $\boldsymbol\Gamma=\frac{\alpha_s}{4\pi}\boldsymbol\Gamma_0+\cdots$，
$\boldsymbol\Gamma'\equiv\partial\boldsymbol\Gamma/\partial\ln\mu=-\gamma_{\rm cusp}\sum_iC_i$。对夸克形状因子（$n=2$），式~\eqref{eq:sudakov-dipole} 重现硬匹配系数的演化 $\dd\ln C_V/\dd\ln\mu$
（\cref{ex:sudakov-dipole-twoparton}）。（本定理只涉及部分子振幅。）
\uses{thm:sudakov-factorization,def:sudakov-cusp,prop:sudakov-webs,prop:sudakov-oneloop}
\end{theorem}

\begin{proofsketch}
(1) \emph{因子化。}把\cref{thm:sudakov-factorization}的论证推广到 $n$ 条宽角外线：领头区域由一个硬子图、每个方向一个喷注子图和一个连接各喷注的软子图组成，
得到振幅的硬—喷注—软因子化\cite{Sen1983}（全阶的振幅层次表述见\cite{FeigeSchwartz2014}，宽角振幅的红外森林公式见\cite{Ma2020}）。
(2) \emph{极点 = 低能矩阵元的紫外极点。}在纯维数正规化中喷注与软矩阵元没有标度，它们的红外极点与紫外极点相消（\cref{ex:sudakov-onshell-hard}的机制），
于是振幅的红外极点由软—共线矩阵元的紫外重整化因子给出；喷注部分是色对角的，软部分是 $n$ 条类光 Wilson 线在原点相遇的关联函数，其重整化在色空间中是乘性的，
web 结构（\cref{prop:sudakov-webs}）保证它可以写成矩阵的指数\cite{Becher:2009cu,Becher:2009qa,DixonMagneaSterman2008}。
(3) \emph{重整化群。}于是 $\mathbf Z$ 满足式~\eqref{eq:sudakov-Zfact}。单圈形式如下：设 $\mathbf Z=1+\frac{\alpha_s}{4\pi}(\frac{\mathbf z_2}{\varepsilon^2}+\frac{\mathbf z_1}{\varepsilon})+O(\alpha_s^2)$，
$D$ 维中 $\frac{\dd\alpha_s}{\dd\ln\mu}=-2\varepsilon\alpha_s+O(\alpha_s^2)$，则
$\frac{\dd\mathbf Z}{\dd\ln\mu}=\frac{\alpha_s}{4\pi}\big[-\frac{2\mathbf z_2}{\varepsilon}-2\mathbf z_1+\frac1\varepsilon\frac{\partial\mathbf z_1}{\partial\ln\mu}\big]$；要求它等于有限的 $-\frac{\alpha_s}{4\pi}\boldsymbol\Gamma_0$，
得 $\mathbf z_1=\frac12\boldsymbol\Gamma_0$、$\mathbf z_2=\frac12\partial\mathbf z_1/\partial\ln\mu=\frac14\boldsymbol\Gamma'_0$，即所述单圈公式。由色守恒，$\sum_{(i,j)}\mathbf T_i\cdot\mathbf T_j=-\sum_i\mathbf T_i^2=-\sum_iC_i$，
故 $\boldsymbol\Gamma'=-\gamma_{\rm cusp}\sum_iC_i$ 是色对角的。
(4) \emph{偶极子形式。}单圈时软关联函数是两两 Wilson 线之间单胶子交换之和，每一对贡献色因子 $\mathbf T_i\cdot\mathbf T_j$ 与类光尖点的反常维度（\cref{def:sudakov-cusp}），由此得到式~\eqref{eq:sudakov-dipole} 的单圈形式；
两圈时软反常维度矩阵与单圈矩阵成正比\cite{Aybat:2006mz}，即偶极子形式在两圈仍成立；Becher 与 Neubert\cite[式~(7)]{Becher:2009cu}（亦见\cite[式~(7)]{Becher:2009qa}）以及 Gardi 与 Magnea\cite{Gardi:2009qi}
猜测它对所有阶成立。两圈以内的显式结构与早先的 Catani 公式\cite{Catani:1998bh}及\cite{Sterman:2002qn}一致。完整论证见上引文献。
\end{proofsketch}

\paragraph{地位与缺口}极点的全阶因子化（式~\eqref{eq:sudakov-Zfact}）是全阶论证（[A]），它继承振幅层次硬—喷注—软因子化的地位（与\cref{thm:sudakov-factorization}同样的两个瓶颈：
减除的 E-S 地位与非阿贝尔 Ward 恒等式）。偶极子形式在单圈与两圈是显式计算的结果（[R] 意义上的定阶内容，两圈见\cite{Aybat:2006mz}）；
作为全阶精确公式它不成立（\cref{prop:sudakov-dipole-violation}）。升级为 $[R]$ 需要振幅因子化的完整证明（参见第~24~章\cref{prop:front-ir-forest}的地位，前向指针）。

\begin{example}[两部分子情形与色守恒]\label{ex:sudakov-dipole-twoparton}
\emph{形状因子。}$n=2$，色单态要求 $\mathbf T_1+\mathbf T_2=0$，故 $\mathbf T_1\cdot\mathbf T_2=-C_F$。两个有序对给出
$\boldsymbol\Gamma=-C_F\gamma_{\rm cusp}\ln\frac{\mu^2}{-s_{12}}+2\gamma^q$。类空形状因子中一入一出，$\sigma_{12}=-1$，$-s_{12}=2p_1\cdot p_2=Q^2$，
\[
\boldsymbol\Gamma=\frac{\alpha_s}{4\pi}\Big[4C_F\ln\frac{Q^2}{\mu^2}-6C_F\Big]=\frac{C_F\alpha_s}{4\pi}(4L-6),\qquad L=\ln\frac{Q^2}{\mu^2}.
\]
另一方面，式~\eqref{eq:sudakov-F1-L} 的有限部分 $C_V=1+\frac{C_F\alpha_s}{4\pi}(-L^2+3L-8+\zeta_2)$ 满足 $\frac{\dd C_V}{\dd\ln\mu}=\frac{C_F\alpha_s}{4\pi}(-2L+3)(-2)=\frac{C_F\alpha_s}{4\pi}(4L-6)$，
两者一致；若把 $(i,j)$ 读作无序对，结果会少一半，与 $C_V$ 不符——这就是式~\eqref{eq:sudakov-dipole} 中“有序对”的含义。% VF-SCET-8, VF-SUD-5
再看极点：$\boldsymbol\Gamma_0=C_F(4L-6)$、$\boldsymbol\Gamma'_0=-8C_F$，单圈公式给出 $\ln\mathbf Z=\frac{C_F\alpha_s}{4\pi}[-\frac2{\varepsilon^2}+\frac{2L-3}\varepsilon]$，正是式~\eqref{eq:sudakov-F1-L} 的极点部分。
类时形状因子中两部分子同为出射，$\sigma_{12}=+1$，$-s_{12}=-s-\ii0$，$\ln\frac{\mu^2}{-s-\ii0}=\ln\frac{\mu^2}{s}+\ii\pi$：$\ln\mathbf Z$ 的极点获得虚部 $\frac{C_F\alpha_s}{4\pi}(-\frac{2\pi\ii}\varepsilon)$，与\cref{prop:sudakov-coulomb}一致。
\emph{三部分子。}对 $q\bar qg$，$\sum_{(i,j)}\mathbf T_i\cdot\mathbf T_j=-(2C_F+C_A)$；SU($N$) 中 $C_A-2C_F=1/N$（脚本 \texttt{s3\_irpoles.py}）。
\uses{thm:sudakov-ir-poles,prop:sudakov-oneloop}
\end{example}

%%% BEGIN fig:sudakov-dipole
\begin{figure}[htbp]
\centering
\begin{tikzpicture}[fd,
  fa soft/.style={gluon,green!45!black},
  fa web/.style={soft blob,line width=0.7pt,pattern=north east lines,pattern color=green!45!black},
  fa over/.style={preaction={draw=white,line width=4.2pt,decorate=false}},
  fa lab/.style={font=\small},
  fa note/.style={font=\footnotesize},
  fa cap/.style={font=\footnotesize},
  fa legs/.pic={\foreach \a/\n in {66/1,22/2,-22/3,-66/4} {
                  \draw[plain] (\a:0.4) -- (\a:2.2) node[font=\small,anchor=\a+180,inner sep=2pt] {$\n$};}
                \node[hard blob,circle,minimum size=0.8cm,inner sep=0pt] at (0,0) {$H$};}]
  % ===================== (a) dipole terms =====================
  \begin{scope}
    \pic at (0,0) {fa legs};
    \draw[fa soft] (66:1.75) to[bend left=38] (22:1.75);
    \draw[fa soft] (-22:1.75) to[bend left=38] (-66:1.75);
    \foreach \a in {66,22,-22,-66} \node[vertex] at (\a:1.75) {};
    \node[fa note,anchor=south west,inner sep=1pt] at (44:2.12) {$\mathbf T_1\cdot\mathbf T_2$};
    \node[fa note,anchor=north west,inner sep=1pt] at (-44:2.12) {$\mathbf T_3\cdot\mathbf T_4$};
    \node[fa note,align=left,anchor=west] at (2.75,0) {对所有\\$i\neq j$ 求和};
    \node[fa cap] at (1.3,-3.05) {(a)\ 两圈以内：偶极子（色关联 $\mathbf T_i\cdot\mathbf T_j$）};
  \end{scope}
  % ===================== (b) three-loop webs (non-dipole) =====================
  % web blob in the free sector between legs 2 and 3; gluons that pass over a leg get a white gap
  \begin{scope}[xshift=7.2cm]
    \pic at (0,0) {fa legs};
    \draw[fa soft,fa over] (1.66,0.27) -- (66:1.5);
    \draw[fa soft] (1.86,0.2) -- (22:2.02);
    \draw[fa soft] (1.86,-0.2) -- (-22:2.02);
    \node[fa web,circle,minimum size=0.62cm,inner sep=0pt] (W) at (1.75,0) {};
    \foreach \p in {(66:1.5),(22:2.02),(-22:2.02)} \node[vertex] at \p {};
    \node[fa note] at (0.9,-2.6) {连接三条外线};
  \end{scope}
  \begin{scope}[xshift=11.6cm]
    \pic at (0,0) {fa legs};
    \draw[fa soft,fa over] (1.66,0.27) -- (66:1.5);
    \draw[fa soft,fa over] (1.66,-0.27) -- (-66:1.5);
    \draw[fa soft] (1.86,0.2) -- (22:2.02);
    \draw[fa soft] (1.86,-0.2) -- (-22:2.02);
    \node[fa web,circle,minimum size=0.62cm,inner sep=0pt] (W) at (1.75,0) {};
    \foreach \p in {(66:1.5),(22:2.02),(-22:2.02),(-66:1.5)} \node[vertex] at \p {};
    \node[fa note] at (0.9,-2.6) {连接四条外线};
  \end{scope}
  \node[fa cap] at (9.45,-3.05) {(b)\ 三圈起：非偶极子 web（阴影）};
  % ===================== anomalous dimension (VF-SUD-5) =====================
  \node[font=\small] at (5.6,-3.9)
    {$\boldsymbol\Gamma=\displaystyle\sum_{(i,j)}\frac{\mathbf T_i\cdot\mathbf T_j}{2}\,\gamma_{\rm cusp}(\alpha_s)\ln\frac{\mu^2}{-s_{ij}}+\sum_i\gamma^i(\alpha_s)$
     \quad（两圈以内成立）};
\end{tikzpicture}
\caption{多腿振幅的红外结构：(a) 两圈以内由偶极子项（色关联 $\mathbf T_i\cdot\mathbf T_j$）给出，反常维度如图下公式（VF-SUD-5，$(i,j)$ 取有序对）；
(b) 三圈起出现连接三条或四条外线的非偶极子 web（偶极子公式作为全阶公式在三圈失效，[V]）。}
\label{fig:sudakov-dipole}
\end{figure}
%%% END fig:sudakov-dipole

\cref{fig:sudakov-dipole}(a) 画出两圈以内的结构：每一对外线之间的软交换只通过 $\mathbf T_i\cdot\mathbf T_j$ 与一个运动学对数进入，就像经典的色偶极子辐射。
这一简单结构的物理来源是：Wilson 线之间的关联在低阶只能通过单胶子交换或其自能修正实现，而这些修正只改变 cusp 反常维度的数值（$\gamma_{\rm cusp}$ 的高阶系数），不产生新的色结构。
从三圈起，一个 web 可以同时连接三条或四条 Wilson 线（\cref{fig:sudakov-dipole}(b)），带来新的色张量与新的运动学依赖。

\begin{proposition}[偶极子公式作为全阶精确公式的失效 \tagV]\label{prop:sudakov-dipole-violation}
式~\eqref{eq:sudakov-dipole} 作为对所有阶成立的公式是错误的：三圈时 $\boldsymbol\Gamma$ 获得非偶极子修正 $\boldsymbol\Delta^{(3)}_n$。
它由两类色结构组成：只含三个不同部分子的项 $\propto C\,f^{abe}f^{cde}\{\mathbf T_i^a,\mathbf T_i^d\}\mathbf T_j^b\mathbf T_k^c$，$C=\zeta_5+2\zeta_2\zeta_3\approx4.99$，对三个有色部分子是一个与运动学无关的常数；
含四个不同部分子的项乘以共形交比的权五单值多重对数函数，对四个及以上有色部分子依赖于运动学。该修正保持类时严格共线因子化\cite[式~(7)、(13)、(14)、(17)]{Almelid:2015jia}。
（本命题只涉及部分子振幅。）
\uses{thm:sudakov-ir-poles}
\end{proposition}

\begin{proofsketch}
（证据：显式计算。）Almelid、Duhr 与 Gardi 计算了连接三条与四条类光 Wilson 线的全部三圈 web，得到完整的三圈软反常维度\cite{Almelid:2015jia}；
Henn 与 Mistlberger 在 $\mathcal N=4$ 超对称 Yang--Mills 理论中计算的三圈四胶子振幅为这一结果提供了独立检验\cite{Henn:2016jdu}。
此前，Caron-Huot 由 Regge 极限中的 BFKL 动力学论证了偶极子猜测从四圈起不成立\cite[第~4.3~节，式~(4.28)]{Caron-Huot:2013fea}。
非偶极子项的结构可以这样理解：三线与四线 web 的色因子中出现两个 $f^{abc}$ 的缩并，它们不能约化为 $\mathbf T_i\cdot\mathbf T_j$；
对三个部分子，运动学依赖只能通过 $\ln(-s_{ij})$ 进入，而重标度不变性（类光 Wilson 线在 $n_i\to\kappa_in_i$ 下的不变性）禁止这种依赖，于是只剩常数；
四个部分子才有重标度不变的共形交比。
\end{proofsketch}

\paragraph{地位与缺口}[V]，证据类型：显式三圈计算（\cite{Almelid:2015jia}，并由\cite{Henn:2016jdu}独立检验）以及四圈的 Regge 论证\cite{Caron-Huot:2013fea}。
被推翻的只是“偶极子形式对所有阶精确”这一陈述；\cref{thm:sudakov-ir-poles}的极点因子化本身（[A]）与两圈以内的偶极子形式不受影响。
三圈修正首先出现在三圈振幅的 $1/\varepsilon$ 极点中；它与类时共线极限相容（严格共线因子化得以保持），这是一个非平凡的检验。
